\chapter{Thiết kế chi tiết}
\label{app:design}

Phụ lục này ghi lại phần chi tiết cài đặt mà Chương~\ref{ch:method} cố ý lược
bỏ: bản đồ module, lược đồ dữ liệu mà mọi giai đoạn trao đổi với nhau, lớp
trừu tượng backend cho phép cùng một pipeline chạy trên mock hoặc trên trình
duyệt thật, và cấu hình chính xác đã dùng cho các run báo cáo ở
Chương~\ref{ch:experiments}.

\section{Bản đồ module}
\label{app:design-modules}

Bảng~\ref{tab:app-modules} liệt kê mọi module của phần cài đặt kèm vai trò của
nó. Phần lõi của pipeline được giữ nhỏ một cách có chủ đích: toàn bộ đường đi
từ một trajectory thô đến một skill được chấp nhận chưa tới 700 dòng, và chính
điều đó làm cho bộ test mock tất định trở nên khả thi.

{\footnotesize
\begin{longtable}{@{}L{4.3cm}rL{8.0cm}@{}}
\caption[Bản đồ module của phần cài đặt HSGE]{Bản đồ module. Số dòng chỉ tính
phần cài đặt, không tính test.}
\label{tab:app-modules}\\
\toprule
\textbf{Module} & \textbf{Dòng} & \textbf{Vai trò} \\
\midrule
\endfirsthead
\toprule
\textbf{Module} & \textbf{Dòng} & \textbf{Vai trò} \\
\midrule
\endhead
\bottomrule
\endfoot
\multicolumn{3}{@{}l}{\emph{Lõi pipeline (\texttt{hsge/})}} \\
\midrule
\texttt{schema.py}      & 129 & Các kiểu dữ liệu dùng chung: \texttt{Step},
\texttt{SubTrajectory}, \texttt{SkillCard}, \texttt{ParamSlot}, \texttt{Task} \\
\texttt{segmentation.py} & 138 & Điểm cắt ứng viên (đổi URL + sụt độ tương tự
axtree), LLM phân xử, gộp đoạn navigation một bước \\
\texttt{clustering.py}  & 100 & Bộ gom cụm phi tham số theo luồng
(ngưỡng cosine / DP-means), không cần biết trước số cụm \\
\texttt{abstraction.py} &  59 & Nâng một cụm chín thành \texttt{SkillCard}
có tham số \\
\texttt{verification.py} &  62 & Ground, thực thi và chỉ chấp nhận skill khi
đạt $\geq k$ lần thành công ở ngữ cảnh held-out \\
\texttt{library.py}     & 105 & \texttt{SkillLibrary}: quan sát $\rightarrow$
gom cụm $\rightarrow$ trừu tượng hóa $\rightarrow$ kiểm chứng $\rightarrow$
chấp nhận; truy hồi và tỉa bớt \\
\texttt{discovery.py}   &  90 & \texttt{Discovery.on\_trajectory}: hook trực
tuyến điều khiển pipeline từ bên trong mọi vòng lặp khám phá \\
\texttt{graph.py}       &  61 & \texttt{HierarchicalGraph}: node skill và cạnh
compose skill$\rightarrow$skill \\
\texttt{operators.py}   & 128 & Operator precondition/effect, planner BFS,
DFS có chặn \texttt{all\_chains} trên đồ thị toán tử \\
\texttt{operator\_induction.py} & 140 & LLM induce precondition/effect từ
skill card cùng các quan sát biên; kiểm tra grounding \\
\texttt{predicate\_registry.py} &  66 & Chuẩn hóa predicate NL tự do thành id
ổn định theo ngưỡng cosine \\
\texttt{curriculum.py}  &  60 & Ghép skill thành chuỗi task dễ$\rightarrow$khó \\
\texttt{transfer.py}    &  53 & Truy hồi, ground lại và kiểm chứng một skill
trên site mới \\
\texttt{llm.py}         & 149 & Giao diện \texttt{LLM.complete\_json} và
\texttt{Embedder.encode}; \texttt{MockLLM}, \texttt{HashingEmbedder} \\
\texttt{backends.py}    & 128 & Backend thật: \texttt{OpenAILLM} (có retry,
tùy chọn vision) và \texttt{STEmbedder} \\
\texttt{env.py}         &  95 & Giao diện \texttt{WebEnv}
(\texttt{reset}/\texttt{step}/\texttt{check}) và các site mock \\
\midrule
\multicolumn{3}{@{}l}{\emph{Runtime WebArena (\texttt{hsge/webarena/})}} \\
\midrule
\texttt{runtime.py}     & 681 & \texttt{WebArenaExplorer}: vòng lặp frontier,
thăng skill$\rightarrow$operator, goal task, lưu trữ \\
\texttt{agents.py}      & 363 & \texttt{PageExplorer} (đề xuất task) và
\texttt{Solver} (vòng lặp hành động), prompt, phát hiện kẹt \\
\texttt{browser.py}     & 143 & \texttt{WebArenaBrowser}: lớp bọc BrowserGym,
quan sát, ảnh chụp màn hình, khai thác href \\
\texttt{auth.py}        & 138 & Luồng đăng nhập theo site và phiên
\texttt{storage\_state} được cache \\
\texttt{predicates.py}  &  61 & Bộ kiểm tra mục tiêu cuối theo URL/DOM, dùng
để chấm thành công của episode \\
\midrule
\multicolumn{3}{@{}l}{\emph{Script ngoại tuyến (thư mục gốc)}} \\
\midrule
\texttt{run\_webarena.py} & 362 & Điểm vào: phân giải cấu hình, nối LLM theo
vai trò, điều phối run \\
\texttt{build\_sft\_dataset.py} & 224 & Episode $\rightarrow$ \texttt{jsonl}
SFT sắp theo curriculum, giữ prompt parity \\
\texttt{judge\_episodes.py} & 108 & LLM kiểm chứng lại các episode tự nhận
thành công \\
\texttt{replay\_segmentation.py} & 131 & Chạy lại segmentation ngoại tuyến trên
các episode đã ghi \\
\end{longtable}}

\section{Lược đồ dữ liệu lõi}
\label{app:design-schema}

Năm dataclass mang toàn bộ dữ liệu đi qua pipeline;
Bảng~\ref{tab:app-schema} liệt kê các trường của chúng. Việc gom chúng vào một
module là điều làm cho các artifact do LLM sinh ra có thể soi xét được, và nhờ
đó mới có được vết kiểm chứng của khâu verification.

\begin{table}[H]
\centering
\caption[Lược đồ dữ liệu lõi]{Các kiểu dữ liệu lõi và trường của chúng.}
\label{tab:app-schema}
\small
\begin{tabular}{@{}L{3.1cm}L{3.6cm}L{6.8cm}@{}}
\toprule
\textbf{Kiểu} & \textbf{Trường} & \textbf{Ý nghĩa} \\
\midrule
\multirow{5}{*}{\texttt{Step}}
 & \texttt{action}         & Chuỗi hành động đã thực thi, vd. \texttt{click('12')} \\
 & \texttt{url}            & URL \emph{sau} hành động \\
 & \texttt{obs\_summary}   & Tóm tắt ngắn bằng ngôn ngữ tự nhiên về kết quả \\
 & \texttt{dom\_signature} & Vân tay cấu trúc rẻ tiền của trang \\
 & \texttt{axtree}         & Axtree đầy đủ (chỉ có ở backend thật) \\
\midrule
\multirow{5}{*}{\texttt{SubTrajectory}}
 & \texttt{steps}        & Lát cắt liên tục của một trajectory \\
 & \texttt{subgoal}      & Mô tả subgoal bằng ngôn ngữ tự nhiên \\
 & \texttt{source\_site} & Tên miền nơi lát cắt được thu \\
 & \texttt{source\_task} & Task sinh ra trajectory cha \\
 & \texttt{success}      & Chỉ các lần thành công mới đủ điều kiện khám phá \\
\midrule
\multirow{4}{*}{\texttt{ParamSlot}}
 & \texttt{name}        & Tên slot được nâng ra, vd. \texttt{username} \\
 & \texttt{slot\_type}  & \texttt{value} $\mid$ \texttt{selector} $\mid$ \texttt{url} \\
 & \texttt{description} & Slot đại diện cho cái gì \\
 & \texttt{example}     & Một giá trị cụ thể đã quan sát được \\
\midrule
\multirow{7}{*}{\texttt{SkillCard}}
 & \texttt{name}, \texttt{description} & Danh tính người đọc hiểu được \\
 & \texttt{precondition}       & Tập khởi tạo $I$, dạng NL kiểm tra được \\
 & \texttt{action\_template}   & Chính sách có tham số $\pi$ \\
 & \texttt{slots}              & Tham số của template \\
 & \texttt{success\_predicate} & Điều kiện kết thúc $\beta$ \\
 & \texttt{stats}              & Số thành viên, số lần kiểm chứng, các site \\
 & \texttt{verified}           & Cờ đã được chấp nhận \\
\midrule
\multirow{3}{*}{\texttt{Task}}
 & \texttt{instruction}  & Mục tiêu ngôn ngữ tự nhiên giao cho solver \\
 & \texttt{skill\_chain} & Các skill id được ghép, theo thứ tự \\
 & \texttt{difficulty}   & Mặc định bằng độ dài chuỗi \\
\bottomrule
\end{tabular}
\end{table}

\noindent
{\sloppy
Sự tương ứng với khung options là trực tiếp và có chủ đích: một
\texttt{SkillCard} chính là một option $\langle I, \pi, \beta \rangle$ với ba
thành phần lần lượt là \texttt{precondition}, \texttt{action\_template} cùng
các \texttt{slots}, và \texttt{success\_predicate}. Còn
\texttt{action\_signature} của một \texttt{SubTrajectory}, chỉ gồm dãy động
từ hành động, là thứ làm cho việc gom cụm không phụ thuộc site: đăng nhập
trên hai site khác nhau cho cùng một signature dù mọi selector và giá trị đều
khác.\par}

\section{Lớp tách backend}
\label{app:design-backends}

Ba giao diện tách pipeline khỏi mọi thứ tốn tiền, cần GPU, hoặc cần site thật:

\begin{itemize}
  \item \texttt{LLM.complete\_json(system, user)} dùng bởi segmentation,
        abstraction, grounding, operator induction và cả hai agent.
        \texttt{MockLLM} trả về đáp án kịch bản; \texttt{OpenAILLM} gọi bất kỳ
        endpoint tương thích OpenAI nào, thử lại tối đa ba lần với backoff mũ
        cho các lỗi tạm thời (rate limit, kết nối, timeout, 5xx - không bao
        giờ thử lại lỗi xác thực hay bad request), và có thể đính kèm ảnh chụp
        màn hình cho model hỗ trợ vision.
  \item \texttt{Embedder.encode(texts)} dùng bởi clustering và predicate
        registry. \texttt{HashingEmbedder} là tất định; \texttt{STEmbedder}
        bọc một model sentence-transformers.
  \item \texttt{WebEnv} với \texttt{reset}, \texttt{step} và \texttt{check}
        pipeline không cần gì hơn từ môi trường. Bản mock cung cấp các site
        kịch bản; \texttt{WebArenaBrowser} điều khiển trình duyệt thật qua
        BrowserGym.
\end{itemize}

\noindent
Nhờ lớp tách hẹp này, toàn bộ pipeline khám phá được kiểm thử bằng các test
tất định khẳng định chính xác số skill và số cạnh compose, mà không cần API
key, GPU hay trình duyệt nào.

\section{Cấu hình model theo vai trò}
\label{app:design-config}

Mỗi thành phần dùng LLM được cấu hình độc lập, nhờ đó phần việc rẻ có thể giao
cho model rẻ trong khi induction và judging dùng model mạnh hơn. Một vai trò
thừa kế mọi trường nó không ghi đè từ khối cấu hình chung; một cờ dòng lệnh
toàn cục ghi đè mọi vai trò cùng lúc. Bảng~\ref{tab:app-roles} ghi lại cách
phân bổ đã dùng cho các run ở Chương~\ref{ch:experiments}.

\begin{table}[H]
\centering
\caption[Phân bổ model theo vai trò]{Phân bổ model theo vai trò. Cột ``Số lần
gọi'' là bậc độ lớn số lần gọi trên mỗi node, chính nó biện minh cho đánh
đổi chi phí/độ mạnh ở từng dòng.}
\label{tab:app-roles}
\footnotesize
\begin{tabular}{@{}L{3.6cm}L{2.8cm}L{2.6cm}L{4.4cm}@{}}
\toprule
\textbf{Vai trò} & \textbf{Model} & \textbf{Số lần gọi} & \textbf{Ghi chú} \\
\midrule
\texttt{page\_explorer} & GPT-4o & một lần/node & Vision + Set-of-Marks,
\texttt{detail=low}, temperature 0,7 để đề xuất đa dạng \\
\texttt{solver} & GPT-4o-mini & tối đa
\texttt{max\_steps}/episode & Vision + Set-of-Marks, temperature 0,0 \\
\texttt{trajectory\_segmenter} & DeepSeek & mỗi trajectory & Chỉ văn bản \\
\texttt{skill\_abstractor} & DeepSeek & mỗi cụm chín & Chỉ văn bản \\
\texttt{skill\_verifier} & DeepSeek & mỗi skill ứng viên & Chỉ văn bản \\
\texttt{operator\_inducer} & DeepSeek & mỗi skill được nhận & Chỉ văn bản \\
\texttt{episode\_judge} & DeepSeek & mỗi lần tự nhận đạt & Chỉ văn bản \\
\bottomrule
\end{tabular}
\end{table}

\section{Tham số chạy}
\label{app:design-params}

Bảng~\ref{tab:app-params} liệt kê các tham số chạy không thuộc diện bí mật.
Các bí mật (API key, mật khẩu tài khoản và URL của site) nằm ngoài kho
mã nguồn, trong một tập tin môi trường. Thứ tự ưu tiên là cờ dòng lệnh, rồi
tập tin cấu hình, rồi giá trị mặc định.

\begin{table}[H]
\centering
\caption[Tham số chạy]{Tham số chạy và giá trị đã dùng.}
\label{tab:app-params}
\small
\begin{tabular}{@{}L{4.4cm}L{1.4cm}L{7.2cm}@{}}
\toprule
\textbf{Tham số} & \textbf{Giá trị} & \textbf{Ý nghĩa} \\
\midrule
\texttt{max\_nodes}     & 500 & Ngân sách node cho một run \\
\texttt{max\_depth}     & 300 & Độ sâu frontier tối đa \\
\texttt{max\_steps}     & 15  & Ngân sách hành động của solver mỗi episode \\
\texttt{goal\_max\_steps} & 25 & Ngân sách bước cho task GOAL và task không tiền tố \\
\texttt{max\_tasks}     & 9   & Số task đề xuất mỗi node \\
\texttt{maturity}       & 2   & Số lần tái xuất hiện để một cụm thành skill \\
\texttt{task\_retries}  & 0   & Số lần thử lại task thất bại (nhân chi phí) \\
\texttt{unprefixed\_sampling} & true & Giải lại mỗi task đã giải được từ trang chủ \\
\texttt{verify\_episodes} & true & Judge chấm lại mọi lần tự tuyên bố thành công \\
\texttt{plan\_subgoals} & true & Phân rã task GOAL thành chuỗi subgoal \\
\texttt{use\_site\_seed} & true & Sinh task GOAL từ seed năng lực của website \\
\texttt{seed\_capability\_repeats} & 3 & Số lần thử lại một năng lực seed trong một run \\
\texttt{goal\_chains\_per\_node} & 3 & Số task GOAL mới thử mỗi node (chuỗi + seed) \\
\texttt{predicate\_sim} & 0,6 & Ngưỡng cosine của predicate registry \\
\texttt{mine\_hrefs}    & true & Khai thác liên kết \texttt{<a href>} vào frontier \\
\texttt{per\_template\_cap} & 3 & Số node nhận vào cho mỗi template URL (tầng 1) \\
\texttt{reached\_template\_cap} & 8 & Trần tương ứng cho tầng 0 (URL đã đi tới) \\
\texttt{max\_hrefs\_per\_page} & 100 & Số liên kết khai thác mỗi trang \\
\texttt{denylist} & - & Chặn URL đăng xuất, phân trang và sắp xếp \\
\bottomrule
\end{tabular}
\end{table}

\chapter{Các prompt}
\label{app:prompts}

Mọi prompt được chép lại nguyên văn, đúng dạng model thực sự nhận được, và giữ
nguyên tiếng Anh vì đó chính là văn bản chạy thật. Chúng được đặt ở đây thay vì
trong Chương~\ref{ch:method} vì cách diễn đạt chính xác của chúng là kết quả
của những sửa chữa hành vi bàn ở Chương~\ref{ch:discussion}: việc phân loại
task có cấu trúc, điều cấm lộ đáp án và điều cấm đề xuất trùng lặp ở
Mục~\ref{app:prompt-explorer} mỗi thứ đều thay thế cho một cách diễn đạt đã bị
một run thật chứng minh là hỏng.

\section{Không gian hành động}
\label{app:prompt-actions}

Được dùng chung bởi prompt của solver và bộ dựng dataset ngoại tuyến, nhờ đó
một mẫu huấn luyện không bao giờ trưng ra tập hành động khác với tập có sẵn lúc
suy luận.

\begin{lstlisting}[style=mystyle,numbers=none,basicstyle=\scriptsize\ttfamily]
Allowed actions (return exactly one): click('bid'), fill('bid','text'),
select_option('bid','opt'), press('bid','key_comb') - e.g. press('12','Enter')
to submit a search box or form when there is no separate submit button,
goto('/path') - a SITE-RELATIVE path on the current site, e.g. goto('/f/books');
never a full http(s):// URL and never another website,
go_back(), new_tab(), tab_focus(index), tab_close(),
scroll(0, 400) - scroll the page down (negative to go up),
hover('bid'), noop(500) - wait 500ms for the page to settle.
If an action failed with a Timeout, the element exists but was not clickable yet
(off-screen, covered, or still loading): scroll(0, 400) or hover it first, or use
a different element - repeating the same click will time out again.
send_msg_to_user('done message'),
report_infeasible('reason') - report_infeasible ONLY when the task is genuinely
impossible on this site, NOT when you are stuck, out of ideas, or cannot find
something yet. If an element is not visible, navigate or look elsewhere first.
\end{lstlisting}

\section{Page explorer}
\label{app:prompt-explorer}

\begin{lstlisting}[style=mystyle,numbers=none,basicstyle=\scriptsize\ttfamily]
PROPOSE concrete, feasible tasks a typical user could do on THIS page.
Return JSON {"tasks": [{"text": "...", "type":
"INFO"|"NAV"|"MOD"|"COMPOSITE"}, ...]}.
Phrase every task the way a real user would state their GOAL or intent - a
natural request, NOT a click-by-click script. Say WHAT they want, not which
buttons to press.
  GOOD: 'How much did I spend on shopping in March?', 'Add a Samsung Galaxy
  S21 to my wish list'.
  BAD (too rigid / step-by-step): 'Click the March filter, then read the
  total'.
Rules:
- Be SPECIFIC using names/values visible on the page (e.g. 'Add the Sony PS5
  to the cart', 'Open the post titled "Best budget GPU"'), NOT generic
  templates like 'add item to cart' or 'search for a product'.
- "type" must be exactly one of: INFO (find/extract a fact), NAV (reach a
  specific page), MOD (change content/state - post, vote, add to cart, edit),
  COMPOSITE (a goal that inherently needs SEVERAL actions to satisfy, stated
  as ONE natural intent - e.g. 'Buy the cheapest in-stock wireless mouse' -
  name the concrete targets from THIS page, but do NOT spell out the
  individual steps). List COMPOSITE tasks FIRST, then MOD, then INFO. Include
  at least 2-3 COMPOSITE tasks whenever the page affords a multi-step goal -
  these are how reusable, chainable skills get discovered - and include a few
  INFO fact-lookup tasks specific to this page. Propose AT MOST ONE NAV task
  (list it LAST), and only to reach an important page that no COMPOSITE/MOD
  task here would visit.
- INFO tasks must name WHAT to find (e.g. 'Find the number of submissions for
  r/relationship_advice'), but must NEVER state the value/answer itself in the
  task text (NOT 'Find the number of submissions for r/relationship_advice
  (5,718 submissions)'). If the answer is already written in the task, it isn't
  a task - whoever solves it must still have to read the page to answer it.
- Do NOT propose a NAV task to a page the user is already on (check the URL
  given below) - e.g. if URL already ends in /wiki, do not propose 'Go to the
  Wiki page'.
- Do NOT re-propose a task you (or an earlier page) already proposed this run
  - see 'ALREADY PROPOSED' below if present. Sitewide nav-menu links
  (header/footer, present on every page) are especially prone to this: once a
  destination has been proposed from one page, do not propose 'go to <that
  same destination>' again from a different page. Prefer content unique to
  THIS page.
- Prefer tasks that make progress or reveal new pages over pure
  search/sort/filter.
- If THIS page has a form with MULTIPLE input fields (e.g. checkout, create a
  listing/post, edit profile, contact/review form), do NOT propose a task that
  fills just one field in isolation ('Enter an email in the email box').
  Propose ONE task that completes the WHOLE form to accomplish its purpose.
  Such a task is naturally COMPOSITE.
- Focus on common things a real user of THIS site would do; skip niche tasks.
- NEVER propose an account/authentication task - log in, log out, sign in,
  sign up, register, create an account, change/reset/forgot password, edit
  email/username, manage security settings, or delete/deactivate/close the
  account. The session is already authenticated and the account is shared
  across exploration runs, so these are useless and risk locking the run out.
\end{lstlisting}

\noindent
Ở chế độ read-only, đoạn ghi chú sau được nối thêm, và mọi task mà bộ phân tích
không phân loại được đều bị loại cùng với mọi task \texttt{MOD}:

\begin{lstlisting}[style=mystyle,numbers=none,basicstyle=\scriptsize\ttfamily]
READ-ONLY MODE: propose ONLY tasks with type INFO or NAV. Do NOT propose any
type MOD / content-modification tasks (no posting, commenting, voting, adding
to cart, or checkout).
\end{lstlisting}

\section{Site seeder}
\label{app:prompt-seeder}

Được gọi đúng một lần cho mỗi lượt chạy, trên trang chủ. Ba ràng buộc
``kiểm chứng được'', ``mọi giá trị đầu vào phải được nêu hoặc suy ra được'' và
``chỉ dùng thứ trang thật sự có'' đều là phản ứng với các mục tiêu seed mà
lượt chạy trước sinh ra rồi không ai chấm được (Mục~\ref{sec:method-seed}).

\begin{lstlisting}[style=mystyle,numbers=none,basicstyle=\scriptsize\ttfamily]
SUMMARIZE this website into a reusable profile for an exploration agent. You
are shown the landing page (its URL and accessibility tree). Infer what KIND of
site this is and the main GOALS a real user comes here to accomplish.
Return JSON: {"description": "<one or two sentences: what this site is and who
uses it>", "capabilities": [{"name": "<short slug, e.g. 'checkout'>", "goal":
"<one natural user intent that exercises this capability>"}, ...]}.
Rules:
- List 4-8 capabilities, ordered from most central to most peripheral.
- Each "goal" must be a concrete, self-contained user intent phrased the way a
  real user would state it (WHAT they want), NOT a click-by-click script.
- Each "goal" must inherently take SEVERAL actions to satisfy - a multi-step
  flow, not a one-shot lookup.
- Each "goal" must END in a CONCRETE, VERIFIABLE state that a checker (or a
  person) could confirm from the resulting page. AVOID open-ended activities
  whose completion is TRIVIAL TO CLAIM: never phrase a goal around 'explore',
  'browse', 'look around', 'view the site/category', or 'find items of
  interest'.
- Every INPUT the goal needs must either be STATED in the goal itself or be
  derivable from the page (e.g. 'the cheapest one', 'the most recent order').
  NEVER leave a value for the agent to invent - it cannot know which one you
  meant and will correctly refuse.
- Ground every capability in what THIS site actually offers - use section/menu
  names visible in the tree. Do NOT invent features you cannot see.
- Prefer capabilities whose completion CHANGES site state or reaches a NEW page
  over pure read-only lookups.
- NEVER include an account/authentication goal.
\end{lstlisting}

\section{Bộ lập kế hoạch subgoal}
\label{app:prompt-planner}

Chỉ được gọi cho các task GOAL, khi \texttt{plan\_subgoals} được bật.

\begin{lstlisting}[style=mystyle,numbers=none,basicstyle=\scriptsize\ttfamily]
DECOMPOSE the user's GOAL into an ordered list of concrete subgoals a web agent
will complete one at a time, in order, to achieve the GOAL on the current site.
You are given the GOAL and the current page (URL + accessibility tree).
Return JSON {"subgoals": ["<subgoal 1>", "<subgoal 2>", ...]}.
Rules:
- 2-5 subgoals. Each is a single, concrete, checkable step stated as an intent,
  NOT a click-by-click script.
- Order them so each subgoal leaves the page in the state the next needs.
- The LAST subgoal must correspond to the GOAL's final end-state (what a grader
  would check), so completing it means the whole GOAL is done.
- Ground subgoals in what THIS page/site actually offers; do not invent steps
  for features you cannot see.
- If the GOAL is already a single concrete step, return it as one subgoal.
\end{lstlisting}

\section{Solver}
\label{app:prompt-solver}

Đây chính là prompt mà model fine-tune được huấn luyện theo; bộ dựng dataset
import trực tiếp hằng số này thay vì chép lại, và đó là điều bảo đảm prompt
parity giữa huấn luyện và suy luận.

\begin{lstlisting}[style=mystyle,numbers=none,basicstyle=\scriptsize\ttfamily]
You are a web agent. Reach the GOAL by choosing the single best next action
from the page accessibility tree (elements shown as [bid] role 'name').
<ACTION SPACE, see B.1>
Think step by step; reflect on your past actions and any error from your last
action before deciding. Output ONE action as JSON
{"thought": "...", "action": "<one action>"}.
Examples: {"thought": "Open the product", "action": "click('12')"} ;
{"thought": "My last action failed because bid 41 does not exist on this page.
Looking at the tree again, the search box is bid 57, so I will type there
instead.", "action": "fill('57','wireless mouse')"} ;
{"thought": "My click on bid 41 timed out, so the element is on the page but
not clickable yet - it is probably below the fold. I will scroll down to bring
it into view before clicking it again.", "action": "scroll(0, 400)"} ;
{"thought": "Goal met", "action": "send_msg_to_user('done: price is $19')"}.
When the `# Error Message from Last Action` section is not (none), your thought
MUST say what went wrong and what you are doing differently - never repeat the
action that just failed.
Call send_msg_to_user ONLY when the goal is fully achieved.
For 'go to page X' goals: once the URL/page has changed to the requested
destination the goal is DONE - call send_msg_to_user immediately; do NOT
click the same link again.
\end{lstlisting}

\section{Induction operator}
\label{app:prompt-operator}

Chỉ dẫn buộc tái dùng nguyên văn một predicate đã biết chính là cơ chế
vocabulary-aware ở Mục~\ref{sec:method-semantic}: chỉ dựa vào độ tương tự
embedding thì không gộp được synonym một cách đáng tin, nên prompt được cung
cấp sẵn bộ từ vựng hiện có và được yêu cầu dùng lại nguyên văn.

\begin{lstlisting}[style=mystyle,numbers=none,basicstyle=\scriptsize\ttfamily]
You convert a reusable web skill into a PLANNING OPERATOR. Given the skill and
the page state before/after it ran on a few sites, state the world conditions
in short, site-agnostic natural language. Preconditions = what must already be
true to apply the skill; effects = what becomes true after.
State preconditions and effects as PAGE STATES - what page or screen the user
is on and what is visible on it (e.g. 'on the product listing page', 'the cart
page is shown') - NOT the mere presence of one element ('a cart link is
present'). Include at least one precondition naming the ENTRY page-state and at
least one effect naming the EXIT page-state, worded so a later skill starting
from that exit state would phrase its own entry precondition identically.
The EXIT page-state MUST DIFFER from the ENTRY page-state. If the skill stays
on the same page, do not repeat the entry phrase - name the same page QUALIFIED
by what the skill established on it, e.g. entry 'on the product edit page' ->
exit 'on the product edit page with a save confirmation'. An effect that merely
restates a precondition is discarded and the whole operator is thrown away, so
a bare repetition wastes the skill. A qualified page-state is still a
page-state: never fall back to an element-level effect ('the description field
is filled') to make it differ - no other skill can phrase that as its entry
precondition. Keep each predicate a terse phrase.
Mark terminal=true only if an effect is a final/transactional state (order
placed, purchase confirmed). A list of KNOWN predicates from previously induced
operators may be provided. If a precondition or effect you want to state is
semantically equivalent to a known predicate, you MUST reuse that predicate's
exact text verbatim instead of rephrasing it. Only coin a new phrase when no
known predicate matches.
Return JSON {"preconditions": ["..."], "effects": ["..."], "terminal": bool}.
\end{lstlisting}

\noindent
Một bước kiểm tra grounding riêng sẽ kiểm lại từng predicate được induce trên
trạng thái trang của \emph{mọi} thành viên trong cụm. Kết quả của nó vừa là
phần đánh giá chất lượng operator ở Chương~\ref{ch:discussion}, vừa là một
cổng chặn thật: effect nào được xác nhận trên dưới một nửa số thành viên sẽ bị
loại, và operator không còn effect nào thì bị từ chối
(Mục~\ref{sec:method-semantic}).

\begin{lstlisting}[style=mystyle,numbers=none,basicstyle=\scriptsize\ttfamily]
You verify PLANNING PREDICATES against real page states. You are given a list
of predicates (each tagged precondition or effect) and, for each occurrence of
the skill, the page's accessibility tree BEFORE and AFTER the skill ran. For a
precondition, check it against the BEFORE tree of every occurrence; for an
effect, check it against the AFTER tree of every occurrence.
Return JSON {"results": [{"predicate": "<exact text>", "type":
"precondition"|"effect", "member_holds": [true|false, ...]}, ...]}, one entry
per predicate, with member_holds having exactly one boolean per occurrence in
the order given.
\end{lstlisting}

\section{Judge cho episode}
\label{app:prompt-judge}

Judge được yêu cầu một cách có chủ đích là chấm trạng thái cuối chứ không chấm
lời tự nhận của agent, đúng thất bại mà nó sinh ra để bắt.

\begin{lstlisting}[style=mystyle,numbers=none,basicstyle=\scriptsize\ttfamily]
You are a strict grader of a web agent's completed task. You are given the GOAL
the agent was told to accomplish, the actions it took before declaring the task
done, and the FINAL page (URL + accessibility tree) where it declared done.
Decide whether the goal was genuinely achieved by the final state - not
whether the agent SAID it was done. Common false positives: the agent gave up
and declared done anyway; it answered an INFO question with a guess not
actually shown on the final page; it navigated to the wrong page.
Return JSON {"verdict": "PASS"|"FAIL", "reason": "<one sentence>"}.
\end{lstlisting}

\chapter{Thống kê run chi tiết}
\label{app:runs}

Phụ lục này ghi lại thống kê đầy đủ của các run khám phá, kể cả những run
ngắn và bị ngắt giữa chừng, để công sức thu thập hiện ra trọn vẹn chứ không
chỉ qua những run đẹp nhất.

Cần phân biệt rõ hai bộ run, vì chúng phục vụ hai mục đích khác nhau và số
liệu của chúng không được cộng gộp ở bất kỳ đâu trong khóa luận:

\begin{itemize}
  \item \textbf{Bốn đợt thu thập chính thức} (Mục~\ref{app:runs-final}) trên
  bốn website của WebArena. Đây là nguồn của mọi con số ở
  Mục~\ref{sec:exp-data} và của toàn bộ tập dữ liệu huấn luyện.
  \item \textbf{Mười hai run phát triển} (Mục~\ref{app:runs-dev}) chỉ trên
  Reddit và GitLab, chạy trong lúc pipeline còn đang được sửa. Chúng không
  đóng góp mẫu huấn luyện nào; giá trị của chúng là các lỗi mà chúng phơi ra,
  và phần lớn cơ chế chốt chặn mô tả ở Chương~\ref{ch:method} ra đời từ đây.
\end{itemize}

\section{Bốn đợt thu thập chính thức}
\label{app:runs-final}

Bảng~\ref{tab:app-runs-final} bổ sung cho Bảng~\ref{tab:runs} của
Mục~\ref{sec:exp-data} những số liệu vận hành mà chương chính không trình
bày: node bị bỏ vì lỗi, số lần khôi phục, và cơ cấu operator bị từ chối.

\begin{table}[H]
\centering
\caption[Số liệu vận hành của bốn đợt thu thập]{Số liệu vận hành của bốn đợt
thu thập chính thức. ``Node lỗi tải'' là node mà trình duyệt không mở được và
cơ chế hỏng mềm đã bỏ qua; ``Reset lỗi'' là số lần đưa trình duyệt về trang
của node trước một task bị thất bại. Hai cột cuối tách số operator bị từ chối
theo lý do.}
\label{tab:app-runs-final}
\small
\begin{tabular}{@{}lrrrrrrrr@{}}
\toprule
& & \multicolumn{2}{c}{\textbf{Sự cố}} & & &
\multicolumn{2}{c}{\textbf{Operator bị loại}} \\
\cmidrule(lr){3-4}\cmidrule(lr){7-8}
\textbf{Website} & \textbf{Node} & \textbf{lỗi tải} & \textbf{reset} &
\textbf{Nối lại} & \textbf{Skill} & \textbf{lặp thừa} & \textbf{vô căn} &
\textbf{Op.} \\
\midrule
Shopping       & 79 & 38 & 166 & 1 & 14 &  1 &  5 &  8 \\
Reddit         & 68 &  0 &   8 & 1 & 16 &  5 &  7 &  4 \\
GitLab         & 63 &  1 &  16 & 1 & 16 &  5 &  3 &  8 \\
Shopping Admin & 48 &  3 &  22 & 2 & 13 &  3 &  5 &  5 \\
\midrule
\textbf{Tổng} & \textbf{258} & \textbf{42} & \textbf{212} & \textbf{5} &
\textbf{59} & \textbf{14} & \textbf{20} & \textbf{25} \\
\bottomrule
\end{tabular}
\end{table}

\noindent
Ba điểm đáng ghi nhận. Thứ nhất, \textbf{cơ chế khôi phục đã được dùng thật
chứ không chỉ được kiểm thử}: sau khoảng hai giờ rưỡi, cả bốn đợt cùng dừng
vì tài khoản API hết hạn mức, và cả năm lần nối lại đều chạy tiếp từ điểm
kiểm tra của node cuối cùng (Mục~\ref{sec:method-impl}). Nếu không có cơ chế
này, toàn bộ phần dữ liệu thu được trong hai giờ rưỡi đầu đã mất.

Thứ hai, \textbf{sự cố trang tập trung gần như trọn vẹn ở Shopping}: 38 trên
42 node lỗi tải và 166 trên 212 lần reset hỏng đều thuộc website này, đều là
timeout điều hướng của trình duyệt. Cơ chế hỏng mềm chuyển từng sự cố thành
một node bị bỏ qua có ghi log, nên không đợt nào bị dừng giữa chừng vì lý do
này.

Thứ ba, \textbf{phần lớn kỹ năng không trở thành operator}: 34 trên 59 kỹ
năng bị chặn lại, 14 vì hiệu ứng lặp lại chính điều kiện tiên quyết và 20 vì
hiệu ứng không được dữ liệu của cụm xác nhận. Mục~\ref{sec:disc-operators}
phân tích thế giằng co giữa hai cổng lọc này.

Số episode ở Bảng~\ref{tab:runs} lấy theo bản tóm tắt cuối cùng mà mỗi đợt
ghi ra (2178 episode). Trên đĩa có 2184 thư mục episode: sáu episode chênh
lệch được ghi ở node đang xử lý dở khi đợt chạy bị dừng, sau lần ghi tóm tắt
cuối cùng, nên chúng không được tính vào bất kỳ thống kê nào.

\section{Mười hai run phát triển}
\label{app:runs-dev}

Bảng~\ref{tab:app-runs} báo cáo các run của giai đoạn phát triển. Run 1-9
khám phá instance Reddit; run 10-12 khám phá GitLab. Run 4-6 là các lượt sinh
task dùng lại đồ thị có sẵn nên không có số node riêng. Trạng thái \emph{bị
ngắt} nghĩa là run được dừng thủ công, không phải là run hỏng; run~10 là run
duy nhất kết thúc bằng lỗi (một timeout điều hướng của trình duyệt mà cơ chế
xử lý lỗi ở mức node đã chuyển thành một node bị bỏ qua trước khi run được
dừng).

\begin{table}[H]
\centering
\caption[Thống kê khám phá theo từng run]{Thống kê khám phá theo từng run. Cột
``Task'' đếm mỗi task đúng một lần, theo kết cục cuối cùng.}
\label{tab:app-runs}
\small
\begin{tabular}{@{}lrrrrrrrr@{}}
\toprule
& & & \multicolumn{3}{c}{\textbf{Task}} & & & \\
\cmidrule(lr){4-6}
\textbf{Run} & \textbf{Node} & \textbf{Ep.} & \textbf{đạt} & \textbf{hỏng} &
\textbf{bất khả} & \textbf{Skill} & \textbf{Op.} & \textbf{Site} \\
\midrule
1  &   1 &   6 &   3 &   3 &  0 &  1 &  1 & Reddit \\
2  &   1 &   8 &   7 &   1 &  0 &  1 &  1 & Reddit \\
3  &  38 & 108 &  71 &  37 &  0 &  5 &  4 & Reddit \\
7  & 128 & 514 & 316 & 198 &  0 & 13 & 13 & Reddit \\
8  &  11 &  43 &  23 &  20 &  0 &  1 &  1 & Reddit \\
9  &  62 & 228 & 182 &  18 & 28 &  4 &  4 & Reddit \\
10 &   7 &  31 &  29 &   1 &  1 &  3 &  3 & GitLab \\
11 &  27 & 134 & 109 &  22 & 15 &  5 &  5 & GitLab \\
12 &  26 & 102 &  87 &  30 &  4 &  4 &  3 & GitLab \\
\midrule
\textbf{Tổng} & \textbf{301} & \textbf{1174} & \textbf{827} & \textbf{330} &
\textbf{48} & \textbf{37} & \textbf{35} & \\
\bottomrule
\end{tabular}
\end{table}

\noindent
Có hai điểm đáng ghi nhận. Thứ nhất, run lớn nhất (run~7, 128 node) ra đời
trước khi có kết cục \emph{bất khả thi}, và đó là lý do tỉ lệ thất bại của nó
cao nhất: những episode mà một run về sau sẽ xếp là bất khả thi thì ở đây bị
tính là thất bại. Thứ hai, hiện tượng trùng lặp node là có thật nhưng ở mức vừa
phải: 29 node trên toàn bộ các run bị bỏ qua vì axtree ban đầu của chúng băm
ra trùng với một node đã thăm, xác nhận rằng chuyện các URL khác nhau hiển thị
cùng nội dung là có thật trên những site này, nhưng không phải là chi phí chi
phối.

\section{Các lượt sinh task của giai đoạn phát triển}
\label{app:runs-gen}

Run 4-6, vẫn thuộc bộ run phát triển, đo riêng khâu đề xuất và kiểm chứng
task trên một đồ thị đã khám phá sẵn. Bảng~\ref{tab:app-gen} tách kết cục của chúng theo loại task. Quy luật
nhất quán ở cả ba lượt: task \texttt{INFO} kiểm chứng đạt ổn định nhất, task
\texttt{MOD} kém nhất, và khoảng một đề xuất trong mười bị bỏ vì đã cũ (trang
mà nó nhắc tới đã thay đổi kể từ lúc đồ thị được ghi.

\begin{table}[H]
\centering
\caption[Sinh và kiểm chứng task theo loại]{Sinh và kiểm chứng task theo loại,
run 4-6.}
\label{tab:app-gen}
\small
\begin{tabular}{@{}lrrrrrrrr@{}}
\toprule
& & \multicolumn{2}{c}{\textbf{INFO}} & \multicolumn{2}{c}{\textbf{NAV}} &
\multicolumn{2}{c}{\textbf{MOD}} & \\
\cmidrule(lr){3-4}\cmidrule(lr){5-6}\cmidrule(lr){7-8}
\textbf{Run} & \textbf{Thử} & \textbf{đạt} & \textbf{hỏng} & \textbf{đạt} &
\textbf{hỏng} & \textbf{đạt} & \textbf{hỏng} & \textbf{Cũ} \\
\midrule
4 & 92 & 43 &  9 & 12 & 8 & 16 & 4 &  8 \\
5 & 89 & 45 &  4 & 18 & 2 & 13 & 7 & 11 \\
6 & 85 & 37 &  8 & 13 & 7 & 15 & 5 & 15 \\
\midrule
\textbf{Tổng} & \textbf{266} & \textbf{125} & \textbf{21} & \textbf{43} &
\textbf{17} & \textbf{44} & \textbf{16} & \textbf{34} \\
\bottomrule
\end{tabular}
\end{table}

\noindent
Gộp lại, task \texttt{INFO} kiểm chứng đạt 86\%, \texttt{NAV} đạt 72\% và
\texttt{MOD} đạt 73\%. Việc \texttt{NAV} không hề dễ hơn \texttt{MOD} mới là
điều phản trực giác: task điều hướng thất bại không phải vì đích đến khó tới,
mà vì agent thường không nhận ra là mình đã tới nơi, đúng hành vi mà mệnh đề
``một khi URL đã đổi thì mục tiêu ĐÃ XONG'' trong prompt của solver
(Mục~\ref{app:prompt-solver}) được thêm vào để sửa.

\section{Goal task được ghép}
\label{app:runs-goal}

Trong suốt \emph{mười hai run phát triển}, không một goal task ghép nào được
thử: bộ đếm chuỗi của cả mười hai bản tóm tắt đều bằng không, và mọi episode
đã ghi đều mang nhãn task mức thấp. Run~7, đồ thị lớn nhất của giai đoạn đó
với mười ba operator, chỉ chứa đúng một cạnh
effect$\rightarrow$precondition, và cạnh ấy chỉ hình thành khi operator thứ
mười ba, cũng là operator cuối cùng, được quy nạp -- thời điểm mà run đã bị
dừng.

\textbf{Bốn đợt thu thập chính thức thì khác: cơ chế ghép chuỗi đã kích hoạt
được trên website sống.} Bảy chuỗi đã được thử, và
Bảng~\ref{tab:app-chains} liệt kê trọn vẹn cả bảy, vì đây là bằng chứng
trực tiếp cho đóng góp lõi của khóa luận nên không nên chỉ báo cáo dưới dạng
một con số tổng.

\begin{table}[H]
\centering
\caption[Bảy chuỗi operator đã thử]{Toàn bộ các chuỗi operator được ghép và
đưa vào thực thi trong bốn đợt thu thập chính thức. Hai lần ``hỏng, 0 bước''
là episode không mở được trang xuất phát, không phải tác nhân thao tác sai.}
\label{tab:app-chains}
\small
\begin{tabular}{@{}lL{7.3cm}rl@{}}
\toprule
\textbf{Website} & \textbf{Chuỗi skill được ghép} & \textbf{Mắt} &
\textbf{Kết cục} \\
\midrule
Shopping & \texttt{navigate\_to\_product\_page} $\rightarrow$
\texttt{add\_variant\_to\_cart} & 2 & đạt, 5 bước \\
Shopping & \texttt{clear\_filter\_via\_clear\_a} $\rightarrow$
\texttt{navigate\_to\_product\_page} & 2 & đạt, 12 bước \\
Shopping & \texttt{navigate\_to\_product\_cate} $\rightarrow$
\texttt{navigate\_to\_product\_page} $\rightarrow$
\texttt{add\_variant\_to\_cart} & 3 & đạt, 7 bước \\
Shopping & \texttt{clear\_filter\_via\_clear\_a} $\rightarrow$
\texttt{navigate\_to\_product\_page} $\rightarrow$
\texttt{add\_variant\_to\_cart} & 3 & hỏng, 0 bước \\
Shopping & \texttt{navigate\_to\_product\_cate} $\rightarrow$
\texttt{navigate\_to\_product\_page} & 2 & hỏng, 0 bước \\
Sh.\ Admin & \texttt{wait\_for\_checkout} $\rightarrow$
\texttt{fill\_checkout\_form} & 2 & đạt, 15 bước \\
GitLab & \texttt{navigate\_to\_fork\_form} $\rightarrow$
\texttt{report\_default\_project\_n} & 2 & hỏng, 1 bước \\
\bottomrule
\end{tabular}
\end{table}

\noindent
Bốn trên bảy chuỗi hoàn thành, trong đó có một chuỗi ba mắt đi trọn đường từ
danh mục sản phẩm tới giỏ hàng trong bảy bước. Điều này trả lời đúng câu hỏi
mà giai đoạn phát triển còn để ngỏ: một website thật \emph{có} cho ra đồ thị
toán tử đủ liên thông để cơ chế ghép chuỗi chạy được.

Nhưng số bảy cũng nói lên giới hạn: 25 operator chỉ sinh ra bảy chuỗi, và ba
trong bốn website đóng góp nhiều nhất một chuỗi. Hệ quả cho tập dữ liệu là
trực tiếp -- chỉ 6 mẫu huấn luyện ở mức độ khó 3 (Bảng~\ref{tab:dataset}),
tức các quỹ đạo dài hạn có mặt trong dữ liệu nhưng quá thưa để chi phối quá
trình huấn luyện. Nguyên nhân nằm ở chỗ một cạnh của đồ thị toán tử đòi hỏi
effect của operator này trùng precondition của operator kia, trong khi sổ
đăng ký mệnh đề chỉ chuẩn hóa bằng độ tương đồng embedding -- thứ mà
Mục~\ref{sec:disc-operators} cho thấy không tách nổi cặp đồng nghĩa thật khỏi
cặp không liên quan ở bất kỳ ngưỡng đơn lẻ nào.

\section{Kết quả đánh giá theo từng task}
\label{app:runs-eval}

\textit{[Sẽ bổ sung khi các lượt đánh giá cho những site còn lại và các lần lặp
thống kê ở Mục~\ref{sec:exp-stats} hoàn tất. Mục này sẽ liệt kê, cho từng task
trong 180 task chính thức của mỗi site, kết cục của bản gốc và bản fine-tune,
để tỉ lệ thành công tổng hợp ở Mục~\ref{sec:exp-main} có thể được soi lại theo
từng task.]}

\chapter{Hướng dẫn tái lập}
\label{app:repro}

Pipeline có thể tái lập trọn vẹn từ một bản checkout sạch. Phụ lục này đưa ra
chuỗi lệnh chính xác; kho mã nguồn đi kèm khóa luận chứa mã, tập tin cấu hình
đã dùng cho các run được báo cáo, và bản hướng dẫn mà mục này cô đọng lại.

\section{Môi trường}
\label{app:repro-env}

Pipeline mock chỉ cần \texttt{numpy} và chạy được mà không cần API key, GPU hay
trình duyệt, đây là cách nên dùng để xem xét hành vi của pipeline trước khi
tiêu tốn bất cứ khoản nào:

\begin{lstlisting}[style=mystyle,numbers=none,basicstyle=\scriptsize\ttfamily]
pip install numpy
python demo.py             # discover -> cluster -> abstract -> verify -> transfer
python explore.py          # exploration loop: frontier -> discover -> compose
python compare_three_way.py  # Flat vs AWM vs HSGE on identical input
python -m pytest hsge/tests/ -q
\end{lstlisting}

\noindent
Pipeline thật cần thêm \texttt{sentence-transformers}, \texttt{openai},
\texttt{browsergym}, \texttt{pyyaml}, \texttt{python-dotenv} và
\texttt{requests}, một instance WebArena tự dựng, cùng thông tin xác thực API
đặt trong một tập tin môi trường không bao giờ được commit.

\section{Thu thập trajectory}
\label{app:repro-collect}

\begin{lstlisting}[style=mystyle,numbers=none,basicstyle=\scriptsize\ttfamily]
# Kiem tra noi day: LLM kich ban va browser gia, khong ton phi, khong can site.
python run_webarena.py -dry-run -out /tmp/wa_dry -max-nodes 2

# Smoke test tren instance dung chung ma khong ghi gi len no.
python run_webarena.py -read-only -max-nodes 2 -max-steps 4

# Run thu thap day du (moi tham so lay tu config.yaml).
python run_webarena.py
\end{lstlisting}

\noindent
Việc đăng nhập cho bốn site có tài khoản chỉ diễn ra một lần và được cache dưới
dạng storage state của trình duyệt; bản cache này là một thông tin phiên còn
sống nên bị loại khỏi quản lý phiên bản. Chỉ chạy một worker: các site WebArena
là trạng thái chia sẻ có thể thay đổi, các run song song sẽ giẫm lên nhau.
Artifact (skill, operator, đồ thị, predicate, kết cục task theo node và một
bản tóm tắt) được ghi sau mỗi node, nên một run bị ngắt vẫn để lại ảnh chụp
nhất quán, và mọi sự kiện đều được ghi song song ra một log JSON lines.

\section{Tinh lọc dataset}
\label{app:repro-curate}

Cả hai lượt kiểm soát chất lượng đều sửa các episode đã ghi tại chỗ, theo kiểu
bổ sung; hãy chạy chúng trên một bản sao nếu cần giữ nguyên artifact gốc.

\begin{lstlisting}[style=mystyle,numbers=none,basicstyle=\scriptsize\ttfamily]
python judge_episodes.py outputs/<run>          # kiem lai cac lan tu nhan dat
python build_sft_dataset.py outputs/<run>/episodes train.jsonl --require-judge
\end{lstlisting}

\noindent
\texttt{--require-judge} loại mọi episode tự nhận thành công nhưng chưa từng
được judge hoặc bị judge đánh trượt, đồng thời vẫn giữ những episode agent báo
là bất khả thi; việc chấm không áp dụng cho chúng và chúng vẫn là mục tiêu
huấn luyện hợp lệ.

\section{Fine-tune, phục vụ và đánh giá}
\label{app:repro-train}

Việc fine-tune dùng LoRA trên model gốc, với phép chia tập được thực hiện theo
trọn trajectory chứ không theo từng bước, để không trajectory nào góp bước cho
cả hai phía. Chat template phải tắt chế độ thinking: dữ liệu thu được không
chứa vết suy luận theo dạng template mong đợi, và để bật sẽ huấn luyện model
sinh ra một phần mà nó chưa từng thấy mục tiêu nào. Sau khi trộn adapter,
model được phục vụ cục bộ và được đánh giá bằng bộ công cụ chính thức của
benchmark chứ không phải một bản cài lại, nhờ vậy số liệu báo cáo mới so sánh
được với các kết quả đã công bố.

\section{Artifact được báo cáo}
\label{app:repro-artifacts}

Sau đây là các artifact cụ thể đứng sau Chương~\ref{ch:experiments}: bốn thư
mục của các đợt thu thập chính thức (Mục~\ref{app:runs-final}) cùng mười hai
thư mục của các run phát triển (Mục~\ref{app:runs-dev}), mỗi thư mục kèm
skill, operator, đồ thị, sổ đăng ký mệnh đề, kết cục theo node, log sự kiện
và điểm kiểm tra nhị phân; các episode đã ghi; tập dữ liệu SFT đã dựng;
adapter LoRA cùng model đã trộn; và kết quả đánh giá của bộ công cụ chính
thức. Mọi con số ở Chương~\ref{ch:experiments} đều lấy từ bốn thư mục đầu;
mười hai thư mục còn lại chỉ được trích dẫn khi nói về các lỗi mà chúng phơi
ra.
